\documentclass[a4paper, 12pt]{report}
\usepackage[latin1]{inputenc}
\usepackage[T1]{fontenc}
\usepackage[english]{babel}

\usepackage[top=3cm, bottom=3cm, left=3cm, right=3cm]{geometry}
\usepackage{setspace}
\usepackage{soul}
\usepackage{color}
\usepackage{graphicx}
\usepackage{amsmath}
\usepackage{amssymb}
\usepackage{rotating}
\usepackage{titling}
\usepackage{hyperref}
\usepackage{pdfpages}
\usepackage{float}
% Algorithms
\usepackage{listings}
\usepackage{algorithm}
\usepackage{algpseudocode}
\usepackage[]{algorithm2e}
\usepackage{moreverb}
% Bibliography
\usepackage[nottoc]{tocbibind}

\lstset{upquote=true,
        columns=flexible,
        keepspaces=true,
        breaklines,
        breakindent=0pt,
        basicstyle=\ttfamily,
        commentstyle=\itshape,
        keywordstyle=\bfseries,
        language=Matlab,
        texcsstyle=*\bfseries}

\postdate{
  \begin{center} 
    \vspace*{3cm}
    \includegraphics[width=10cm]{logo_Poly.jpg}
  \end{center}
}

\newcommand{\ora}[1]{\overrightarrow{#1}}
\newcommand{\R}{\mathbb{R}}

\algnewcommand\algorithmicforeach{\textbf{for each}}
\algdef{S}[FOR]{ForEach}[1]{\algorithmicforeach\ #1\ \algorithmicdo}

\begin{document}

\def\labelitemi{$\bullet$}

\title{{\Large \textbf{Artificial Intelligence: Probabilistic \& learning techniques\\ } \\Chris \bsc{Pal} \\ Project (INF8225)}}
\author{Farnoush \bsc{Farhadi} \& Juliette \bsc{Tibayrenc}} 
\date{21 Apr. 2016}

\maketitle

\setcounter{tocdepth}{5}

\tableofcontents

\newpage

\chapter*{Introduction}
\addcontentsline{toc}{chapter}{Introduction}

Our AI course has allowed us to discover Markov Decision Processes, which we decided to further explore
in our project.
Looking for an application for them, we stumbled upon Mastronarde \& Van der Schaar's article,
`Fast Reinforcement Learning for Energy-Efficient Wireless Communications'\cite{Mastronarde-2011},
and we chose it as the basis for our work.

The article presents the two researchers' work in developing an algorithm that would enable systems
to adopt the best policy to quickly send data packets in the most efficient way possible,
an important goal to achieve since the majority of today's surfing is done on smartphones
whose battery does not last longer than a day (at best).

In this report, we will first introduce the article, with a quick review of the state of the art
at the time it was written (2011), then summarize it.
We will continue with a presentation of some theoretical elements, highlighting the differences between
what we saw in class and the processes on which the article and our implementations are based,
and the way the reward matrix used is calculated (a process that does not appear in the article).
Finally we will talk about our own implementation of classic Q-learning and one of its variants,
speedy Q-learning, and compare the results we obtain with the results obtained with the authors'algorithm,
plus the results we get by using value iteration and policy iteration.

\chapter*{1. Reference article}
\addcontentsline{toc}{chapter}{1. Reference article}

\section*{1.1 Literature review}
\addcontentsline{toc}{section}{1.1 Literature review}

Many studies have been performed on energy consumption of Wireless Sensor Networks (WSN), mainly
focused on the transfer and the routing of the data. A low level model for lithium-ion batteries and a high level model for the time-varying load were proposed in \cite{Salodkar-2008} that require a predetermined workload profile which is not considered as a priori known. This workload evolves as a random process. In these problems, the battery behavior is treated as a stochastic process whose parameters are extracted from the electrochemical characteristics of the simulated battery. Numerous stochastic models have been proposed in the literature which uses a discrete-time Markovian chain model. In fact, this approach is a Markovian chain of the battery states of charge with forward and backward sweeps regarding to the normal discharge and recovery effect processes, respectively.
Some state-of-the-art researches divided the problem of energy efficient wireless communications into two categories: physical (PHY) layer-centric solutions and adaptive modulation and coding (AMC) and system-centric solutions such as dynamic power management (DPM). These techniques can be used to trade-off delay and energy to increase the lifetime of battery-operated mobile devices by proposing a meticulous and unified framework for simultaneously utilizing
physical and system level techniques to minimize power consumption subjecting to delay constraints when we have stochastic and unknown traffic and channel conditions. They formulate the power management problem as a Markov decision process and solve it using reinforcement learning method \cite{Mastronarde-2011}.

\section*{1.2 Article summary}
\addcontentsline{toc}{section}{1.2 Article summary}

In this project, we consider the problem of energy-efficient point-to-point
transmission of delay-sensitive data in dynamic environments based on Mastronarde \& Van der Schaar's paper. In such environments there are time-varying channel conditions (e.g., fading channel) and dynamic traffic loads in which the main concern is the reliable delivery of data to the receiver within a tolerable delay. 
The research for addressing this problem can be roughly divided into two categories: 
physical (PHY) layer-centric solutions such as power-control and adaptive modulation and coding (AMC) 
and system-centric solutions such as dynamic power management \cite{Mastronarde-2011}.
Although these techniques are different significantly, they
can be applied to trade-off delay and energy to enhance the lifetime of battery consumption devices.
An excessive amount of PHY-centric methods focus on optimal single-user power-control to minimize
the transmission power subject to queuing delay constraints. 
The main disadvantage of such methods is that they typically require statistical knowledge of the underlying dynamics like the channel state and traffic distributions, which is not always available. 
PHY-centric methods are effective for minimizing transmission power. 
But, a considerable amount of power are wasted in such systems for staying the wireless
card on to be ready to transmit. 
To address this problem, System-level solutions are applied that  rely on DPM, which enables system components to be put into low-power states when they are not required. 
In this paper authors consider the structure of the problem to the DPM and exhibit suboptimal learning performance. They adapt the power-control, AMC, and DPM policies according to the current and future traffic and channel conditions. 
Furthermore, they provide a unified framework for jointly optimizing system-level and PHY-centric power management for delay-sensitive wireless at the first time.
They suggest a novel decomposition of the value iteration and Reinforcement Learning algorithms based on dividing the system’s dynamics into a priori known and unknown elements. Using these partial system information requires to be learned less than when using conventional Reinforcement Learning algorithms and avoids the need for action exploration, which leads to speed up the adaptation and performance of existing Reinforcement Learning methods. 
The method does not need any priori information about the traffic arrival and channel statistics
to determine the jointly optimal physical-layer and system-level power management strategies; 
It also applies partial information about the system in a way that less information required to be learned than when using conventional reinforcement learning algorithms; and by using this approach, we don't need for action exploration, which severely restricts the adaptation speed and run-time performance of conventional reinforcement learning methods.
\\ 
We try to implement their proposed framework for simultaneously utilizing physical and system-level techniques to minimize energy consumption, under delay constraints, in the presence of stochastic and unknown traffic and channel conditions. 
The problem is formulated as a Markov decision process and solve it using value iteration and policy iteration methods. We aim to find the optimal policy (directly or indirectly by optimizing the
value function). We also implement another solution based on Q-learning consists of dynamics and reward function initially unknown to learn the best policy from history.
We compare the MDP results with that of obtained by Q-learning methods.

\chapter*{2. Theoretical elements}
\addcontentsline{toc}{chapter}{2. Theoretical elements}

\section*{2.1 Value iterating \& policy iterating vs. Q-learning}
\addcontentsline{toc}{section}{2.1 Value iterating \& policy iterating vs. Q-learning}

In our project, we compare the performance of several algorithms against that of the authors.
These algorithms have to be evaluated along several criteria, among which the elements that need to be known
for them to work, the speed with which they are executed (partly dependent on the efficency of the implementation and the language used)
and the number of iterations their internal loop is executed for.

These algorithms are value iteration, policy iteration, classic Q-learning and speedy Q-learning.

\subsection*{2.1.1 Value iteration}
\addcontentsline{toc}{subsection}{2.1.1 Value iteration}

One method to find an optimal policy is to find the optimal value function by determining a simple iterative algorithm called value iteration.
The value iteration method is referenced by Bellman equation for MDP problems which is obtained simply by turning the Bellman optimality equation into an update rule.
If there exist $n$ possible states, then we have $n$ Bellman equations, one per each state.
The $n$ equations contain $n$ unknowns which are the utilities of theses states. 

We aim to solve these simultaneous equations to find the utility values.
These equations are non-linear, since the max operator is not a linear one.
Hence, it can be determined by a simple iterative algorithm that can be shown as follows:

\[V_{i+1}(s) = R(s) + \gamma \mathrm{max_a} \sum_s P(s,a,s') V_{i}(s')\]

where $i$ is the iteration number.
Value iteration starts at $i = 0$ and then iterates, repeatedly computing $V_{i+1}$ for all states s, until V converges.
Value iteration formally requires an infinite number of iteration to converge to $V^*$.
In practice, it stops once the value function changes by a small amount in a sweep which satisfies the Bellman equations.
This process is called a value update or Bellman update/backup. 
The algorithm for value iteration is as follows (from Artificial Intelligence, 3rd edition, French version \cite{Russell-2010}):

  
\begin{algorithm}[H]
\KwData{$P(s,a,s')$ /* a transition matrix */, $a$ /* actions */, $R(s)$ /* reward signal */, $\gamma$ /* discount*/, $\epsilon$ /* the maximum error allowed in the utility of any state */}
\KwResult{$V(s)$ /* Value function */}
    Initialize:\\  
    $V(s) \gets \texttt{arbitrarily}$, for all $s \in \mathcal S^+$\\
    Repeat\\
    $V \gets V'$ \\
    \ForEach {$\texttt{state  }$ $s \in \mathcal S^+$}{
    $V^{'}(s) \gets R(s) + \gamma \mathrm{max_a} \sum_{s'} P(s,a,s') V(s')$
    $\Delta \gets \max(\Delta, V(s) - V'(s)|)$}
    Until {$\Delta < \epsilon$}\\
    return $V(s)$
\caption{Iterative Value Algorithm}
\end{algorithm}

Note that the value iteration backup is identical to the policy evaluation except that it requires the max operator to be taken over all actions.

\subsection*{2.1.2 Policy iteration}
\addcontentsline{toc}{subsection}{2.1.2 Policy iteration}

In policy iteration algorithm, the agent manipulates the policy directly, instead of finding it indirectly via the optimal value function. 
The policy often becomes exactly optimal long before the utility estimation have converged to their true values. This phenomena leads us to another way of finding optimal policies, called policy iteration. Policy iteration picks an initial action arbitrarily by taking rewards on states as
their utilities and computing a policy according to the maximum expected
utility. Then it iteratively performs two steps: firstly, it determine the values through
computing the utility of each state given the current action, and then policy improvement by updating the current policy if possible. 
To detect when the algorithm has converged, it should only change the policy if the new action for some state improves the expected value; that is, it should stabilize the policy.
The algorithm for policy iteration is as follows:

\begin{algorithm}[H]
\KwData{$P(s,a,s')$ /* a transition matrix */, $a$ /* actions */, $R(s)$ /* reward signal */, $\gamma$ /* discount*/, $\epsilon$ /* the maximum error allowed in the utility of any state */}
\KwResult{$V(s)$ /* Value function */}
Initialize \\
$V(s) \gets \texttt{arbitrarily}$, for all $s \in \mathcal S^+$\\
$\pi(s) \gets \texttt{arbitrarily}$, for all $s \in \mathcal S^+$\\
$\mathrm{Policy Evaluation:}$\\
$V \gets V^{'}$\\
\ForEach {$\texttt{state  }$ $s \in \mathcal S^+$}{
$V^{'}(s) \gets R(s) + \gamma \mathrm{max_a} \sum_{s'} P(s,a,s') V(s')$\\
$\Delta \gets \max(\Delta, V(s) - V^{'}(s)|)$
}
until {$\Delta < \epsilon$}\\
$\mathrm{Policy Improvement:}$\\
$\mathrm{stable == True}$\\
\ForEach {$\texttt{state  }$ $s \in \mathcal S^+$}{
$a \gets \pi(s)$\\
$\pi(s) \gets \mathrm{argmax}_a \sum_{s'} P(s,a,s')V(s')$\\
\If{$ a \neq \pi(s)$}{
$\mathrm{stable == False}$
}}
\caption{Iterative Policy Algorithm}
\end{algorithm}

At the end of algorithm, we test the stable variable; if it is 'True', algorithm terminates, otherwise, control jumps to the Policy Evaluation and continues in an iterative manner. 
Policy iteration often converges in few iterations. The policy improvement theorem guarantees that these policies works better than the original random policy and achieves to the goal states in the minimum number of steps.

\subsection*{2.1.3 Classic Q-learning}
\addcontentsline{toc}{subsection}{2.1.3 Classic Q-learning}

Q-learning is part of a second category of algorithms that do not presuppose the full knowledge of the system.
Contrary to value iteration and policy iteration, the environment does not need to be known for the algorithm to work.
The algorithm presented by the authors of our reference article is similar in its requirements, though it goes further and
actually uses the knowledge we do have about the environment, which makes it more efficient.

Q-learning helps us find the optimal policy by using the interactions between the system and the environment.
In it, we compute an estimate of the cumulative discounted reward for each action executed in each state.

The algorithm for one step is as follows (from Artificial Intelligence, 3rd edition, French version\cite{Russell-2010}:

\begin{algorithm}[H]
  \KwData{$s'$ /* current state */, $r'$ /* reward signal */}
  \KwResult{$Q$ /* action values table indexed by states and actions */, $T$ /* frequency of state-action couples */, $s,a,r$ /*previous state, action and reward */}
  \If{s final}{
    $Q(s,Empty) = r'$
  }
  \If{$s\neq0$}{
    $T(s,a)++$\\
    $Q(s,a)+=\alpha\cdot T(s,a)(r+\gamma\max_{a'}Q(s',a')-Q(s,a))$
  }
  $s,a,r=s',\argmax_{a'}f(Q(s',a'),T(s',a')),r'$\\
  return $a$
\caption{Q-learning algorithm for one step}
\end{algorithm}

\subsection*{2.1.4 Speedy Q-learning}
\addcontentsline{toc}{subsection}{2.1.4 Speedy Q-learning}

Speedy Q-learning is developed in the article `Speedy Q-learning' (Ghavamzadeh et al. \cite{ghavamzadeh2011speedy}).
It is, as its name indicates, a variant of Q-learning that converges faster than the form we presented in the previous subsection.
We use the synchronous version, as presented in the paper, since we do have a generative model.
This is, however, a restriction on the actual effectiveness of applying this algorithm that we will have to keep in mind later when comparing results.

The algorithm is presented in \cite{ghavamzadeh2011speedy} as follows:

\begin{algorithm}[H]
  \KwData{$Q_0$ /* Initial value for Q */, $\gamma$ /* discount factor */, $T$ number of iterations,$R$ reward matrix}
  \KwResult{$Q_T$}
  $Q_{-1} = Q_0$;\\
  \For{$k=0:1:T-1$}{
    $\alpha_k=\frac{1}{k+1}$;\\
    \Foreach{$(x,a)$}{
      Generate next state sample $y_k$\\
      $\Tau_k Q_{k-1}(x,a)=R(x,a)+\gamma\max_{a\in A}{Q_{k-1}(y_k,a)}$;\\
      $\Tau_k Q_{k}(x,a)=R(x,a)+\gamma\max_{a\in A}{Q_{k}(y_k,a)}$;// Empirical Bellman operator\\
      $Q_{k+1}(x,a)=Q_k(x,a)+\alpha_k(\Tau_kQ_{k-1}(x,a)-Q_k(x,a))+(1-\alpha_k)(\Tau_kQ_{k}(x,a)-\Tau_kQ_{k-1}(x,a))$;//SQL update rule
    }
  }
  return $Q_T$
\caption{Speedy Q-learning}
\end{algorithm}

We notice that in the update rule, the second term is made more important by its learning step $1 - \alpha_k$.
The 

\section*{2.2 Reward matrix}
\addcontentsline{toc}{section}{2.2 Reward matrix}

\chapter*{3. Comparing algorithms}
\addcontentsline{toc}{chapter}{3. Comparing algorithms}

\section*{3.1 Classic Q-learning}
\addcontentsline{toc}{section}{3.1 Classic Q-learning}

\section*{3.2 Speedy Q-learning}
\addcontentsline{toc}{section}{3.2 Speedy Q-learning}

\section*{3.3 Additional remarks}
\addcontentsline{toc}{section}{3.3 Additional remarks}

\chapter*{Further work \& conclusion}
\addcontentsline{toc}{chapter}{Further work and conclusion}

\bibliographystyle{ieeetr}
\bibliography{report}


\end{document}
